\section{载流导线在磁场中的受力}

\subsection{安培力}
在磁场中的载流导线也会受到磁场力的作用，这种力称为\uwave{安培力}。而我们知道载流导线是由大量运动电荷构成，而上一节中我们已经知道运动电荷会受到洛伦兹力，因此宏观上的安培力实际上就是微观上的洛伦兹力集体作用的结果，现在我们就来推导安培力的计算公式。

若我们以$\dif{\vb*{F}}$表述电流源受到的安培力，而以$\vb*{F}$表示每个运动电荷的洛伦兹力
\begin{Equation}
    \dif{\vb*{F}}=\vb*{F}nS\dif{l}=\vb*{F}n\dif{\vb*{V}}=\vb*{F}\dif{n}
\end{Equation}
这样就建立了电流元的安培力$\dif{\vb*{F}}$和运动电荷的洛伦兹力$\vb*{F}$的关系。

根据\Refx{公式}{洛伦兹力}和\Refx{公式}{电流元与运动电荷}
\begin{Equation}
    \dif{\vb*{F}}=\vb*{F}\dif{n}=qv\times\vb*{B}\dif{n}=I\dif{\vb*{l}}\times\vb*{B}
\end{Equation}
即
\begin{BoxEquation}[安培力]
    \dif{\vb*{F}}=I\dif{\vb*{l}}\times\vb*{B}
\end{BoxEquation}
\Refx{公式}{安培力}是对电流元而言的，对于完整的载流导线$L$，应以积分形式表示
\begin{BoxEquation}[安培力的积分形式]
    \vb*{F}=\Ilnt{I\dif{\vb*{l}}\times\vb*{B}}
\end{BoxEquation}
现在我们来考虑一些特殊情况，若载流导线是直导线，而磁场又是匀强磁场时，此时的磁感强度$\vb*{B}$可以作为常量提出积分外，而$\dif{\vb*{l}}$的积分即为由导线一端指向另一端的常矢量$\vb*{l}$
\begin{BoxEquation}[匀强磁场中的载流直导线]
    \vb*{F}=I\vb*{l}\times B
\end{BoxEquation}
值得注意的是，对于匀强磁场中任意形状的载流导线，由于$\dif{\vb*{l}}$的积分的结果总是由导线一端指向另一端的矢量，与导线的具体路径无关（因为$f(x,y)=1$是路径无关的），因此
\begin{BoxPrinciple}[匀强磁场中的安培力性质]
    在均匀磁场中，任意弯曲载流导线的受力等同于从其起点到终点通有相等电流的直导线。

    在均匀磁场中，任意闭合载流导线受到的安培力为零。
\end{BoxPrinciple}
特别注意的是以上结论仅适用于匀强磁场，否则$\vb*{B}$无法作为常矢量提出积分外。

\subsection{平行电流间的安培力}
载流直导线可以在磁场中受到安培力，载流直导线也可以同时产生磁场。

在本小节，我们就来计算一组平行无限长载流直导线中的电流间的相互作用。

由于电流产生的磁场总是沿着电流的右手螺旋方向（且在平行电流的任意直线上的磁场是相等的），而电流受到的安培力是右手四指由电流转向磁场时的拇指指向，因此不难得出\footnote{这里的绘图不是很准确，直导线的磁场是距离反比衰减的，图中以匀强磁场绘之，不过这并没有什么妨碍。}

\begin{Figure}{电流间的相互作用力}
    \subfigure[同向电流相吸]
    {
        \inputfigure[scale=0.6]{Chapter10E_Figure01}
    }\qquad
    \subfigure[异向电流相斥]
    {
        \inputfigure[scale=0.6]{Chapter10E_Figure02}
    }\vspace{-10pt}
\end{Figure}
\Refx{图}{电流间的相互作用力}的结论可以表述如下
\begin{itemize}
    \item 两根电流方向\textbf{相同}的无限长平行电流间的磁场力表现为\textbf{吸引力}（同向相吸）。
    \item 两根电流方向\textbf{不同}的无限长平行电流间的磁场力表现为\textbf{排斥力}（异向相斥）。
\end{itemize}
电流中“同向相吸异向相斥”，电荷间“同性相斥异性相吸”，应注意两者是不同的。

现在我们来进一步的做一些定量计算，根据\Refx{公式}{载流无限长直导线的磁场}
\begin{equation*}
    B_1=\frac{\mu_0I_1}{2\pi d}\qquad
    B_2=\frac{\mu_0I_2}{2\pi d}
\end{equation*}
这里的$\vb*{B}_1,\vb*{B}_2$分别是电流$I_1,I_2$在空间中产生的磁场。

现在我们再来计算载流导线受到的磁场力的大小，虽然由另一直导线产生的磁场并不是匀强的，但是在与之平行的直导线所在直线上的各点处的磁感强度都是相等的，因此\Refx{公式}{匀强磁场中的载流直导线}仍然适用。然而其指出直导线受到的磁场力与它的长度成正比，那么对于这里的无限长直导线，两者间的磁场力就是无穷大，这是没有意义的\footnote{其实这也是合理的，就像密度一定的两根无限长的物体间的万有引力是无穷大一样。}，但是我们仍然可以研究无限长直导线中的某一段在另外一根无限长直导线的磁场中受到的安培力，这将是个有限值，在$I_1,I_2$上任取长度为$l_1,l_2$两端，其受力$\vb*{F}_1,\vb*{F}_2$分别为
\begin{equation*}
    \begin{cases}
        F_1=B_2I_1l_1=\dfrac{\mu_0}{2\pi}\dfrac{I_1I_2}{d}l_1\\[5mm]
        F_2=B_1I_2l_2=\dfrac{\mu_0}{2\pi}\dfrac{I_1I_2}{d}l_2
    \end{cases}
\end{equation*}
由此可见，对于一组平行无限长载流直导线，相同长度的导线总是受到相同的力，无论我们研究的这段导线是具体属于其中哪一根载流直导线，因此上式可以统一表述为
\begin{BoxEquation}[平行载流导线间的作用力]
    F=\frac{\mu_0}{2\pi}\frac{I_1I_2}{d}l
\end{BoxEquation}
\Refx{公式}{平行载流导线间的作用力}是意义是，对于两根载有电流$I_1,I_2$的平行无限长直导线，若两根导线相距为$d$，那么对于其中任何一根导线上长度为$l$的导线段，它将会受到大小为$F$的磁场力，若$I_1,I_2$方向相同则$F$为引力，若$I_1,I_2$方向反向则为斥力。

\Refx{公式}{平行载流导线间的作用力}可以定义国际单位制中的电学基本单位安培
\begin{BoxGeneral}[安培][定义]
    真空中两根相距$1$\si{m}的无限长平行圆直导线中通以等量恒定电流时，

    若每$1$\si{m}长度上的导线受到的作用力为$2\times 10^{-7}$\si{N}，则此时每根导线中的电流为$1$\si{A}。
\end{BoxGeneral}
这里定义了安培后，真空磁导率的数值就被确定了
\begin{Equation}
    \mu_0=4\pi\times 10^{-7}\si{N\cdot A^{-2}}
\end{Equation}
这也就解释了为何真空磁导率的数值是如此整齐了，因为其直接基于一个基本单位的定义。

由于安培的大小被确定后库仑的大小也就确定了，因此根据\Refx{定律}{库仑定律}，如果我们现在按照这个定义取两个$1$库仑的点电荷，使其相距$1$米，那么两者间的库仑力是可以通过实验测定的，这样真空电容率的数值也可以被确定了，而实验表明其值为
\begin{Equation}
    \varepsilon_0=8.85\times 10^{-12}\si{C^2\cdot N^{-1}\cdot m^{-2}}
\end{Equation}
这也就表明，真空磁导率和真空电容率之间是存在耦合关系的，两者不是相互独立的。

这里可以设想，如果我们将\Refx{定义}{安培}中同等条件下的电流改定义为$k$\si{A}，那么上面的点电荷的电荷量也对应的变为$k$\si{C}，如果记此时的真空电容率和真空磁导率为$\varepsilon_0',\mu_0'$
\begin{equation*}
    F_e=\frac{1}{4\pi\varepsilon_0'}\frac{q_1'q_2'}{r^2}\qquad 
    F_m=\frac{\mu_0'}{2\pi}\frac{I_1'I_2'}{d}l
\end{equation*}
而在原先定义下有（注意$F_e,F_m$是不会随着安培和库仑的定义而变的）
\begin{equation*}
    F_e=\frac{1}{4\pi\varepsilon_0}\frac{q_1q_2}{r^2}\qquad 
    F_m=\frac{\mu_0}{2\pi}\frac{I_1I_2}{d}l
\end{equation*}
其中新定义中的$q_1',q_2',I_1',I_2'$是原定义中的$q_1,q_2,I_2,I_2$数值上的$k$倍，因而
\begin{equation*}
    \varepsilon_0'=a^2\varepsilon_0\qquad
    \mu_0'=\frac{\mu_0}{a^2}
\end{equation*}
关注到
\begin{equation*}
    \varepsilon_0'\mu_0'=\varepsilon_0\mu_0
\end{equation*}
这也就是说，当我们修改安培的定义时，尽管真空电容率和真空磁导率的数值会发生变化，但是两者的乘积$\varepsilon_0\mu_0$是一个定值，我们后面还将看到$\varepsilon_0\mu_0$实际与光速$c$有关。

